\section{Random access to the complete evolving field}\label{sec:access}
The formulas above describe more than the time of an event. They give a succinct
representation of the entire field at any requested time. A finite arithmetic
segment $[x,y,d_x,d_y,n]$ denotes
$\{(x+id_x,y+id_y):0\le i<n\}$. Combining segments by XOR means that a cell
is black exactly when an odd number of the segments contain it.

\begin{theorem}[Exact complete-state random access]\label{thm:access}
For all nonnegative binary integers $l,r,T$, the complete state $F^T(X(l,r))$
is computable in polynomial bit time in $\log(l+r+T+2)$. Its field is an XOR
of at most 20,000 finite segments. Every nonsingleton has $d_x=2$ and
$d_y\in\{-2,0,2\}$; coordinates and lengths have $O(\log(l+r+T+2))$ bits.
The representation gives the exact pose, next turn, cell membership and total
black mass, and supports exact equality of two complete fields.
With $M$ explicitly supplied initial toggles inside the certified unread
half-strips, the bound is $20,000+M$ and time is polynomial in total input length.
\end{theorem}
Here $F$ is the ordinary one-update map. The input is the compact pair of counter
lengths, and the output is a succinct complete field. We use the original
coordinate frame of Section~\ref{sec:basin} throughout.

\subsection{Powered words retain every write}
Suppose a certified word $w$ restores heading, has displacement $d$, and has
phase read positions $q_i$ relative to its entry head $p$. For $n$ full copies,
\begin{equation}\label{eq:powered-field}
 p_{\rm new}=p+nd,\qquad
 B_{\rm new}=B\xor\mathop{\bigtriangleup}_{i<|w|}
           \{p+q_i+jd:0\le j<n\}.
\end{equation}
Every update flips exactly its read cell, so this identity preserves all
self-intersections and earlier colors with the correct parity. A partial copy
is handled by quotient and remainder and its finite list of reads. The counter
and recovery words have displacement $(\pm2,0)$; the standard tail has diagonal
displacement with components $\pm2$. Reversing a negatively directed finite
segment normalizes $d_x$ to 2. Singletons can use direction $(2,0)$.

\subsection{Locating the requested time}
Let $m=\min(l,r)$. Before first exit, the largest completed round is
\[
 j=\min\left(m,\left\lfloor
       \frac{\operatorname{isqrt}(\lfloor T/2\rfloor+49)-7}{2}
       \right\rfloor\right).
\]
Indeed $t_j=8j^2+56j=2(2j+7)^2-98$. The integer square root therefore inverts
the exact round inequality without approximation. Write the complete checkpoint
$\J((5)\R(j,l-j))\R(j,r-j)$ directly and apply
\eqref{eq:powered-field} to at most one remaining pair of excursions.
With $P=\word{RLLR}$ and the kernel words $K,E$ of
Section~\ref{sec:five}, the right success is $P^jKP^j$ and its terminal call
is $P^jE$. The left words are the letterwise R/L swaps of
$\word{LR}P^jKP^j\word{RL}$ and $\word{LR}P^jE$, respectively.
Their execution follows from the proved reflection/complement identities.

For $T$ after first exit, start directly with the complete exit formula
\eqref{basin:eq:exit}. Scan the fixed list of recovery blocks, subtracting
lengths arithmetically until the requested block is found. Its partial execution
uses \eqref{eq:powered-field}. The nine infinite parameter domains are covered
by the exact affine verifiers. The 59 exceptional clipped sources have bounded
parameters and recovery words of length at most 17,106, so their words may be
computed by a fixed finite execution. Once $T\ge\tau$, divide $T-\tau$ by 104.
The resulting finite diagonal segments describe every written cell of the
whole growing wake, together with the remainder of the current period.

\subsection{Restoring the original unbounded counter}
Let $(l_c,r_c)$ be the clipped pair in Corollary~\ref{basin:cor:clip}, and set
$s=(2(l-l_c),0)$. If $B_0$ is the complete initial source, put
\[
 D=B_0(l,r)\xor(B_0(l_c,r_c)+s).
\]
Every point of $D$ is unread for the entire common future by the proved
protection guards. First-disagreement induction gives
\[
 p_T(l,r)=p_T(l_c,r_c)+s,\qquad
 B_T(l,r)=(B_T(l_c,r_c)+s)\xor D.
\]
Both source fields in $D$ are constant-sized arithmetic expressions. Thus the
extra counter remains present in the output. Explicit toggles in the certified
unread half-strips append singleton terms; their guard is checked at input.

\subsection{Uniform size and exact set operations}
Each source, checkpoint or exit field has fewer than 15 blocks of width at most
four and at most five black cells per column. A repeated block has width two,
so fewer than 300 segments suffice for each complete such field. The remaining
counter round has phase budget below 150; the largest generic recovery program
has budget 4,074, plus at most 19 remainder reads. A standard tail adds at most
$104+103=207$ segments. The largest exceptional branch is bounded by
\[
 300+17,106+207+2\cdot300=18,213<20,000.
\]
The last two terms restore the source difference. Cancellation only reduces
the count. Coordinates have magnitude $O(l+r+T+1)$ by the source formulas and
unit-step motion. The algorithm uses a bounded number of integer operations,
including one square root, on numbers of logarithmic bit length. This proves
the representation and bit-complexity claims.

Membership is the parity of divisibility and interval tests on the segments.
For total mass, group segments by direction $d_y$, $x$ parity $a$, and intercept
$y-d_y(x-a)/2$. XOR-sweep their integer endpoints, obtaining disjoint intervals
on each line. Distinct parallel lines are disjoint. Two nonparallel lines
intersect in at most one point; exact integer determinants and interval tests
identify every such intersection. Group coincident points and collect all
incident segment identifiers, so triple intersections are counted together.
If $m(p)$ is the number of retained segments meeting at an intersection, then
\[
 |B|=\sum_i |S_i|-\sum_p\bigl(m(p)-(m(p)\bmod2)\bigr).
\]
Ordinary points occur once; this correction replaces each crossing multiplicity
by its XOR parity. There are polynomially many endpoints and pair intersections
in the number of terms, regardless of numeric segment length. Applying the
same mass calculation to the XOR of two fields decides equality. Membership
at the returned pose gives the next turn. These operations complete the proof
of Theorem~\ref{thm:access}.\qed
